% @econ-contract section
% id: CONCLUSION_ANSWER
% role: Answer the question at its proved scope: the mechanism and answer first; who decides accommodation, who decides intensity; what the reach results add; the persistence and welfare limits.
% question_answered: What did the paper establish, and where does it stop?
% required_inputs: RESULTS_HORIZONS, WELFARE_BOUNDARY, DISCUSSION_INTERPRETATION.
% output_handoff: None.
% claim_ceiling: The master plan's ceiling; no policy prescription.
% exclusions: A theorem list; new claims.
% exact_ref: tex-label:sec:conclusion
% support_status: VERIFIED
% status: HUMAN_CONFIRMED
% human_confirmed_at: 2026-10-01
% decision_ref: suggestions/2026-10-01-manuscript-rebuild/01_STRUCTURE_PROPOSAL.md#author-decisions-recorded
% confirmation_scope: Section role and claim ceiling.
% reopen_if: A principal result or the claim ceiling changes.
% amended_at: 2026-10-12 (opens on the dual-basin illustration and states its limits; decision_ref: suggestions/2026-10-10-jet-submission/DECISIONS.md#framing-2026-10-12); 2026-10-13 (argument plan: mechanism and answer first, then the dual-basin illustration with its limits, the reach results, welfare and lock-in, and a closing 'so what / now what' paragraph; decision_ref: suggestions/2026-10-10-jet-submission/exposition-2026-10-13/ARGUMENT_PLAN.md); 2026-10-15 (the dual-basin illustration dropped from the main text for now, on the author's word); 2026-10-19 (coauthor wording review: wording of the coauthor's line-by-line revision applied where not wrong; sender/target in lower case; non-compliance for defiance; decision_ref: suggestions/2026-10-10-jet-submission/DECISIONS.md#coauthor-wording-review-2026-10-19); 2026-10-20 (humanizer pass: dashes replaced by commas, parentheses or colons; plain wording; meaning unchanged; decision_ref: suggestions/2026-10-10-jet-submission/DECISIONS.md#humanizer-pass-2026-10-20)
% @end-econ-contract
\section{Conclusion}
\label{sec:conclusion}

% @econ-contract paragraph
% id: CONCLUSION_ROLES
% role: Open with the mechanism and the answer on the class (the myopic benchmark ends sanctions, forward-looking governments sustain them; the Target's patience decides relief against defiance), then the Sender's role and the same-policy comparison with its hypothesis (repair C4). Moved first (argument plan, 2026-10-13); reordered on the layout decisions of 2026-10-16.
% depends_on: CONCLUSION_ANSWER
% handoff: The reach results.
% exact_ref: tex-label:prop:myopic-benchmark
% exact_ref: tex-label:thm:regime-separation
% exact_ref: tex-label:thm:four-cells
% exact_ref: tex-label:thm:dominance
% support_status: VERIFIED
% status: PROPOSED
% @end-econ-contract
\central{Economic sanctions can persist because they strengthen the domestic political
interests that support their continuation.} Without foresight they unravel: in the
myopic benchmark, the target accommodates and the political state erodes over time. In
the equilibria constructed over the class of economies, when both governments value the
political state and are forward-looking, the target maintains non-compliance, sanctions
persist indefinitely, and the political state converges to a positive steady-state
level. The target's patience determines whether relief or non-compliance prevails;
whether the sender is forward-looking shapes the compliance threshold for relief or the
intensity of sanctions under non-compliance. These comparisons hold all non-discount
primitives fixed. For a fixed target policy, a non-decreasing sender continuation value
weakly increases the largest of the optimal sanctions relative to the myopic best response. Under
interiority, the constructed equilibria exhibit a strictly positive foresight premium
relative to the myopic best response, with a closed-form lower bound.

% @econ-contract paragraph
% id: CONCLUSION_REACH
% role: Summarize the reach results: patience, like the Target's other parameters, can change whether full defiance is an equilibrium; result names in place of condition lists (argument plan, 2026-10-13).
% depends_on: CONCLUSION_ROLES
% handoff: The welfare findings and lock-in.
% exact_ref: tex-label:prop:security
% exact_ref: tex-label:prop:collapse-threshold
% support_status: VERIFIED
% status: PROPOSED
% @end-econ-contract
\central{Beyond the constructed equilibria, patience, like all other target
parameters, shapes whether full non-compliance is supportable as an equilibrium.} When
the sender's response and the target's value under full non-compliance do not decrease
with the state, greater target patience preserves the full-non-compliance equilibrium.
Under the strong entrenchment condition, the target guarantees itself a minmax payoff
against every pure sanction strategy, bounding its payoff in every equilibrium from
below; and across the class, every equilibrium with a non-decreasing sanction rule
exhibits full non-compliance. In the special case of \Cref{prop:collapse-threshold}, a
full-non-compliance equilibrium exists if and only if the target's discount factor is at
least a threshold value.

% @econ-contract paragraph
% id: CONCLUSION_LIMITS
% role: State the welfare findings, then lock-in and persistence without hysteresis, naming prop:lock-in instead of listing its conditions (argument plan, 2026-10-13).
% depends_on: CONCLUSION_REACH
% handoff: So what and now what.
% exact_ref: tex-label:sec:welfare
% exact_ref: tex-label:prop:lock-in
% support_status: VERIFIED
% status: PROPOSED
% @end-econ-contract
\central{Persistent sanctions can be socially costly, and the game alone does not
determine their welfare ranking.} Across the class, two interest-group representations
of the same government payoffs rank the constructed equilibria in opposite orders, and
under compliance neutrality one makes persistent play socially optimal. For an example
economy with one explicit interest-group polity, equilibrium sanctions remain at their
maximum level and generate welfare losses.

The political state can also lock in sanctions. Where an equilibrium and a high-state
interval satisfy the conditions of \Cref{prop:lock-in}, a state entering the interval
remains there under equilibrium play and sanctions never decline. When full sanctions
raise every state below the basin, a sufficiently long temporary shock can produce a
durable outcome. Persistence can also arise without hysteresis. Where the explicit
Lipschitz constant derived for the constructed equilibrium's law of motion is strictly
below one, a temporary shock has no permanent asymptotic effect.

% @econ-contract paragraph
% id: CONCLUSION_SO_WHAT
% role: Say what the results mean (persistence need not reveal failed bargaining or misperception; with information complete, the levers act on the political state and the governments' incentives, not on beliefs) and what would discipline the channel (measurement, sec:discussion:scope), with no prescription; close on the paper's answer.
% depends_on: CONCLUSION_LIMITS
% handoff: None.
% origin: argument plan item 10 (2026-10-13); the closing sentence moved here from CONCLUSION_LIMITS.
% exact_ref: tex-label:sec:discussion:scope
% exact_ref: tex-label:sec:discussion:misperception
% support_status: VERIFIED
% status: PROPOSED
% @end-econ-contract
\central{Persistent sanctions therefore need not reflect failed bargaining or
misperception.} Within the model, they can be the optimal choice of governments that
value the political state that sanctions themselves have built. With complete
information, the levers that end sanctions operate through the political state and
government incentives, not through beliefs (\Cref{sec:discussion:misperception} of the
Online Appendix). Whether this channel operates in a given episode is an empirical
question: measurements that separate the political state from aggregate income would
discipline estimates of its transition and of the payoff terms governing basin entry
(\Cref{sec:discussion:scope}). \central{Optimizing governments can thus perpetuate and
lock in sanctions through the political state that sanctions generate; patience and
other economic features determine whether this occurs.}
